\documentclass[11pt,a4paper]{article}
\usepackage[utf8]{inputenc}
\usepackage[T1]{fontenc}
\usepackage[portuguese]{babel}
\usepackage{xcolor}
\usepackage{hyperref}
\usepackage{geometry}
\usepackage{tcolorbox}

\geometry{margin=2.5cm}

\definecolor{primary}{RGB}{39, 174, 96}
\definecolor{secondary}{RGB}{44, 62, 80}
\definecolor{codebg}{RGB}{240, 240, 240}

\hypersetup{colorlinks=true, linkcolor=primary, urlcolor=primary}

\title{\color{secondary}\Huge\textbf{Solução: Exercício 4.6 - O projeto v2.0}}
\author{\color{primary}DevOps com Kubernetes -- 2026}
\date{}

\begin{document}
\maketitle

\section*{\color{primary}Guia de Solução}
\begin{tcolorbox}[colback=green!5!white,colframe=green!60!black,title=Resolução Passo a Passo]
### Passo a Passo

1. \textbf{Identificar o Problema}: Ao escalar o serviço \textit{broadcaster} (que assina os tópicos do NATS), várias réplicas receberão a mesma mensagem, gerando envios duplicados.
2. \textbf{Uso de Queue Groups no NATS}: Na biblioteca ou cliente NATS utilizado no seu código, modifique a inscrição (\textit{subscription}) para pertencer a um \texttt{Queue Group}.
3. \textbf{Exemplo de Código}: Se for Node.js, em vez de \texttt{nc.subscribe('todo_topic', (msg) => {...})}, altere para \texttt{nc.subscribe('todo_topic', { queue: 'broadcaster-workers' }, (msg) => {...})}.
4. \textbf{Implantação e Teste}: Atualize o código, gere a imagem, mude as réplicas do Deployment do Broadcaster para 6 (\texttt{replicas: 6}). Teste criar uma tarefa. Graças ao Queue Group, o NATS balanceará o tráfego, garantindo que apenas 1 das 6 réplicas processe a mensagem.
\end{tcolorbox}

\end{document}
